\documentclass[12pt, letterpaper]{article}

% ─── Paquetes ───────────────────────────────────────────────────────────────
\usepackage[utf8]{inputenc}
\usepackage[T1]{fontenc}
\usepackage[spanish]{babel}
\usepackage[margin=2.5cm]{geometry}
\usepackage{lmodern}
\usepackage{microtype}
\usepackage{setspace}
\usepackage{parskip}
\usepackage{titlesec}
\usepackage{titletoc}
\usepackage{tocloft}
\usepackage{booktabs}
\usepackage{tabularx}
\usepackage{array}
\usepackage{longtable}
\usepackage{xcolor}
\usepackage{colortbl}
\usepackage{fancyhdr}
\usepackage{graphicx}
\usepackage{enumitem}
\usepackage{hyperref}
\usepackage{float}
\usepackage{subcaption}

% ─── Colores ─────────────────────────────────────────────────────────────────
\definecolor{primary}{HTML}{185FA5}
\definecolor{primarylight}{HTML}{E6F1FB}
\definecolor{primarymid}{HTML}{B5D4F4}
\definecolor{rowalt}{HTML}{F7FAFD}
\definecolor{headerbg}{HTML}{185FA5}
\definecolor{headerfg}{HTML}{FFFFFF}
\definecolor{tablegray}{HTML}{F1EFE8}
\definecolor{bordergray}{HTML}{CCCCCC}
\definecolor{textmuted}{HTML}{5F5E5A}

% ─── Tipografía y espaciado ───────────────────────────────────────────────────
\setstretch{1.25}
\setlength{\parindent}{0pt}
\setlength{\parskip}{8pt}

% ─── Encabezados de sección ──────────────────────────────────────────────────
\titleformat{\section}
  {\large\bfseries\color{primary}}
  {\thesection.}{0.6em}{}[\vspace{2pt}\color{primarymid}\hrule\vspace{4pt}]

\titleformat{\subsection}
  {\normalsize\bfseries\color{primary}}
  {\thesubsection}{0.6em}{}

\titleformat{\subsubsection}
  {\normalsize\bfseries\color{textmuted}}
  {\thesubsubsection}{0.6em}{}

\titlespacing{\section}{0pt}{18pt}{8pt}
\titlespacing{\subsection}{0pt}{12pt}{4pt}

% ─── Encabezado y pie de página ──────────────────────────────────────────────
\pagestyle{fancy}
\fancyhf{}
\fancyhead[L]{\small\color{textmuted}\textit{Silk Road — Especificación de Requerimientos de Software}}
\fancyhead[R]{\small\color{textmuted}2026}
\fancyfoot[C]{\small\color{textmuted}\thepage}
\renewcommand{\headrulewidth}{0.4pt}
\renewcommand{\headrule}{\color{primarymid}\hrule width\headwidth height\headrulewidth}

% ─── Listas ──────────────────────────────────────────────────────────────────
\setlist[itemize]{noitemsep, topsep=4pt, leftmargin=1.4em}
\setlist[enumerate]{noitemsep, topsep=4pt, leftmargin=1.6em}

% ─── Comandos auxiliares ─────────────────────────────────────────────────────

% Fila de encabezado para tablas
\newcommand{\tableheadrow}[1]{%
  \rowcolor{headerbg}\textcolor{headerfg}{\textbf{#1}}%
}

% Etiqueta de requerimiento (RF / RNF / CU)
\newcommand{\reqid}[1]{%
  \textcolor{white}{\colorbox{primary}{\texttt{\small #1}}}%
}

% ─── Hipervínculos ───────────────────────────────────────────────────────────
\hypersetup{
  colorlinks=true,
  linkcolor=primary,
  urlcolor=primary,
  citecolor=primary,
  pdftitle={ERS - Silk Road},
  pdfauthor={Juancarlos Monsalve}
}

% ════════════════════════════════════════════════════════════════════════════
\begin{document}
% ════════════════════════════════════════════════════════════════════════════

\begin{center}
    {\normalsize
    Instituto Tecnológico de Costa Rica \\
    Centro Académico de Alajuela \\
     IC-5821. Requerimientos de software
    }

    \vspace{1.5cm}

    \includegraphics[width=0.75\textwidth]{Logo.png}

    \vspace{1.5cm}

    {\bfseries
    Proyecto - I Parte: \\
    }

    \vspace{1cm}

    {\bfseries
    Estudiantes: \\
    }
    José Pablo Barrantes Hernández - 2025063829  \\
    José David Quirós – 2025065801  \\
    Carlos Andrés Zúñiga – 2025065194 \\
    Juan Carlos Monsalve – 2023141521 \\

    \vspace{1cm}

    {\bfseries
    Profesor:
    }

    Bryan Tomas Hernandez Sibaja

    \vspace{1cm}

    {\bfseries
    Fecha de entrega:
    }

    {\textbf{22/04/2026}}

    \vspace{2cm}

    I Semestre, 2026
\end{center}

\newpage


% ─── Tabla de contenidos ─────────────────────────────────────────────────────
\tableofcontents
\newpage

% ════════════════════════════════════════════════════════════════════════════
\section{Introducción}
% ════════════════════════════════════════════════════════════════════════════

El presente documento tiene como propósito definir de manera detallada los requerimientos del sistema \textbf{SilkRoad}, el cual será desarrollado utilizando el framework de React + Python. Este documento servirá como base para el diseño, desarrollo y validación del sistema.

% ─────────────────────────────────────────────────────────────────────────────
\subsection{Alcance}

El sistema permitirá a los usuarios finales de Silk Road contar con las siguientes funcionalidades:

\begin{itemize}
  \item Autenticación de usuarios (registro, inicio y cierre de sesión)
  \item Exploración y búsqueda de productos
  \item Gestión de carrito de compras y pedidos
  \item Creación y administración de tiendas por parte de vendedores
  \item Publicación y gestión de productos e inventario
  \item Panel de administración para moderación de la plataforma
\end{itemize}

% ─────────────────────────────────────────────────────────────────────────────

% ─────────────────────────────────────────────────────────────────────────────
\subsection{Visión general del documento}

Este documento está organizado en las siguientes secciones:

\begin{enumerate}
  \item \textbf{Introducción}: propósito, alcance, términos y referencias.
  \item \textbf{Visión del Producto}: declaración de visión en formato caja.
  \item \textbf{Descripción General}: contexto, funciones, usuarios y dependencias.
  \item \textbf{Requerimientos Funcionales}: comportamientos esperados del sistema.
  \item \textbf{Requerimientos No Funcionales}: atributos de calidad del sistema.
  \item \textbf{Casos de Uso}: descripción textual de interacciones actor--sistema.
  \item \textbf{Restricciones}: limitaciones técnicas y de proyecto.
  \item \textbf{Modelo de Datos}: estructura de entidades y relaciones.
  \item \textbf{Mapa de Navegación}: estructura de pantallas de la aplicación.
\end{enumerate}

\newpage
% ════════════════════════════════════════════════════════════════════════════
\section{Visión del Producto}
% ════════════════════════════════════════════════════════════════════════════

La siguiente tabla presenta la declaración de visión del producto SilkRoad.

\vspace{8pt}
\begin{tabularx}{\linewidth}{>{\bfseries\color{primary}}p{3.2cm}X}
  \toprule
  \rowcolor{headerbg}
  \multicolumn{2}{l}{\textcolor{white}{\textbf{Silk Road — Declaración de Visión}}} \\
  \midrule
  Para
    & Compradores que buscan productos de distintos vendedores en un solo lugar, y
      vendedores que desean crear y administrar su propia tienda en línea sin
      infraestructura técnica propia. \\
  \rowcolor{rowalt}
  Quienes
    & Necesitan una plataforma centralizada para explorar, comprar y vender productos
      de forma organizada y segura, sin depender de múltiples sitios independientes. \\
  El producto
    & Silk Road es una aplicación web de tipo marketplace multi-tienda
      desarrollada con Django. \\
  \rowcolor{rowalt}
  Que
    & Permite a compradores explorar productos, gestionar un carrito y realizar
      pedidos; a vendedores crear tiendas, publicar productos y administrar su
      inventario; y a administradores moderar y supervisar toda la plataforma. \\
  A diferencia de
    & Soluciones complejas como Amazon o MercadoLibre, que requieren gran
      infraestructura, integraciones de pago avanzadas y procesos de \emph{onboarding}
      extensos para vendedores. \\
  \rowcolor{rowalt}
  Nuestro producto
    & Ofrece una experiencia simplificada y accesible, diseñada para entornos
      educativos y pequeños comercios, con una arquitectura limpia basada en Django
      que facilita su comprensión, mantenimiento y extensión por parte del equipo
      de desarrollo. \\
  \bottomrule
\end{tabularx}

% ════════════════════════════════════════════════════════════════════════════
\section{Descripción General}
% ════════════════════════════════════════════════════════════════════════════

Silk Road es una aplicación web cliente-servidor desarrollada en Django que permite la interacción de los usuarios mediante un navegador web estándar. La plataforma centraliza la oferta de múltiples vendedores independientes en un único punto de acceso para los compradores.

\newpage
% ─────────────────────────────────────────────────────────────────────────────
\subsection{Funciones del sistema}

El sistema proporcionará las siguientes funciones principales:

\begin{itemize}
  \item Registro e inicio de sesión de usuarios
  \item Exploración y búsqueda de productos por categoría o nombre
  \item Gestión de carrito de compras
  \item Creación y procesamiento de pedidos
  \item Creación y administración de tiendas por vendedores
  \item Publicación, edición y eliminación de productos
  \item Control de inventario por parte del vendedor
  \item Panel de administración para gestión de la plataforma
\end{itemize}

% ─────────────────────────────────────────────────────────────────────────────
\subsection{Características de los usuarios}

\begin{tabularx}{\linewidth}{>{\bfseries}lX}
  \toprule
  \rowcolor{primarylight}
  Rol & Descripción \\
  \midrule
  Cliente (comprador)
    & Puede registrarse, navegar el catálogo, buscar productos, añadir artículos al
      carrito, realizar pedidos y consultar su historial de compras. \\
  \rowcolor{rowalt}
  Vendedor
    & Puede crear una tienda dentro de la plataforma, publicar y administrar
      productos, gestionar su inventario y visualizar los pedidos recibidos. \\
  Administrador
    & Gestiona usuarios, aprueba o suspende tiendas, modera productos y revisa
      información general del sistema. \\
  \bottomrule
\end{tabularx}

% ─────────────────────────────────────────────────────────────────────────────
\subsection{Suposiciones y dependencias}

\begin{itemize}
  \item Los usuarios tendrán acceso a internet y a un navegador web moderno.
  \item El equipo utilizará Django como framework principal de desarrollo.
  \item Los requerimientos no cambiarán significativamente durante el ciclo de
        desarrollo del proyecto.
  \item No se integrará un sistema de pago real en esta primera versión.
\end{itemize}

\newpage
% ════════════════════════════════════════════════════════════════════════════
\section{Requerimientos Funcionales}
% ════════════════════════════════════════════════════════════════════════════

\newcolumntype{L}{>{\raggedright\arraybackslash}X}

% ── RF-01 ────────────────────────────────────────────────────────────────────
\vspace{6pt}
\begin{tabularx}{\linewidth}{>{\bfseries\color{primary}}p{3.5cm}L}
  \toprule
  \rowcolor{headerbg}
  \multicolumn{2}{l}{\textcolor{white}{\textbf{RF-01 — Registro de usuario}}} \\
  \midrule
  Descripción   & El sistema debe permitir a nuevos usuarios crear una cuenta
                  proporcionando nombre, correo electrónico y contraseña. \\
  \rowcolor{rowalt}
  Prioridad     & Alta \\
  Rol afectado  & Cliente, Vendedor \\
  \bottomrule
\end{tabularx}

% ── RF-02 ────────────────────────────────────────────────────────────────────
\vspace{6pt}
\begin{tabularx}{\linewidth}{>{\bfseries\color{primary}}p{3.5cm}L}
  \toprule
  \rowcolor{headerbg}
  \multicolumn{2}{l}{\textcolor{white}{\textbf{RF-02 — Inicio de sesión}}} \\
  \midrule
  Descripción   & El sistema debe permitir a los usuarios autenticarse mediante
                  correo electrónico y contraseña. \\
  \rowcolor{rowalt}
  Prioridad     & Alta \\
  Rol afectado  & Cliente, Vendedor, Administrador \\
  \bottomrule
\end{tabularx}

% ── RF-03 ────────────────────────────────────────────────────────────────────
\vspace{6pt}
\begin{tabularx}{\linewidth}{>{\bfseries\color{primary}}p{3.5cm}L}
  \toprule
  \rowcolor{headerbg}
  \multicolumn{2}{l}{\textcolor{white}{\textbf{RF-03 — Exploración de productos}}} \\
  \midrule
  Descripción   & El sistema debe mostrar un catálogo de productos con nombre,
                  precio, imagen y tienda de origen. \\
  \rowcolor{rowalt}
  Prioridad     & Alta \\
  Rol afectado  & Cliente \\
  \bottomrule
\end{tabularx}

% ── RF-04 ────────────────────────────────────────────────────────────────────
\vspace{6pt}
\begin{tabularx}{\linewidth}{>{\bfseries\color{primary}}p{3.5cm}L}
  \toprule
  \rowcolor{headerbg}
  \multicolumn{2}{l}{\textcolor{white}{\textbf{RF-04 — Búsqueda de productos}}} \\
  \midrule
  Descripción   & El sistema debe permitir buscar productos por nombre o categoría. \\
  \rowcolor{rowalt}
  Prioridad     & Media \\
  Rol afectado  & Cliente \\
  \bottomrule
\end{tabularx}

% ── RF-05 ────────────────────────────────────────────────────────────────────
\vspace{6pt}
\begin{tabularx}{\linewidth}{>{\bfseries\color{primary}}p{3.5cm}L}
  \toprule
  \rowcolor{headerbg}
  \multicolumn{2}{l}{\textcolor{white}{\textbf{RF-05 — Carrito de compras}}} \\
  \midrule
  Descripción   & El sistema debe permitir al cliente agregar, modificar y eliminar
                  productos de su carrito antes de confirmar un pedido. \\
  \rowcolor{rowalt}
  Prioridad     & Alta \\
  Rol afectado  & Cliente \\
  \bottomrule
\end{tabularx}

% ── RF-06 ────────────────────────────────────────────────────────────────────
\vspace{6pt}
\begin{tabularx}{\linewidth}{>{\bfseries\color{primary}}p{3.5cm}L}
  \toprule
  \rowcolor{headerbg}
  \multicolumn{2}{l}{\textcolor{white}{\textbf{RF-06 — Creación de pedidos}}} \\
  \midrule
  Descripción   & El sistema debe permitir al cliente confirmar su carrito y generar
                  un pedido. El sistema registrará el pedido y notificará al vendedor. \\
  \rowcolor{rowalt}
  Prioridad     & Alta \\
  Rol afectado  & Cliente \\
  \bottomrule
\end{tabularx}

% ── RF-07 ────────────────────────────────────────────────────────────────────
\vspace{6pt}
\begin{tabularx}{\linewidth}{>{\bfseries\color{primary}}p{3.5cm}L}
  \toprule
  \rowcolor{headerbg}
  \multicolumn{2}{l}{\textcolor{white}{\textbf{RF-07 — Creación de tienda}}} \\
  \midrule
  Descripción   & El sistema debe permitir a un usuario registrado crear una tienda
                  indicando nombre, descripción y categoría principal. \\
  \rowcolor{rowalt}
  Prioridad     & Alta \\
  Rol afectado  & Vendedor \\
  \bottomrule
\end{tabularx}

% ── RF-08 ────────────────────────────────────────────────────────────────────
\vspace{6pt}
\begin{tabularx}{\linewidth}{>{\bfseries\color{primary}}p{3.5cm}L}
  \toprule
  \rowcolor{headerbg}
  \multicolumn{2}{l}{\textcolor{white}{\textbf{RF-08 — Gestión de productos}}} \\
  \midrule
  Descripción   & El sistema debe permitir al vendedor crear, editar y eliminar
                  productos de su tienda. Cada producto tendrá: nombre, descripción,
                  precio, cantidad en inventario e imagen. \\
  \rowcolor{rowalt}
  Prioridad     & Alta \\
  Rol afectado  & Vendedor \\
  \bottomrule
\end{tabularx}

% ── RF-09 ────────────────────────────────────────────────────────────────────
\vspace{6pt}
\begin{tabularx}{\linewidth}{>{\bfseries\color{primary}}p{3.5cm}L}
  \toprule
  \rowcolor{headerbg}
  \multicolumn{2}{l}{\textcolor{white}{\textbf{RF-09 — Visualización de pedidos recibidos}}} \\
  \midrule
  Descripción   & El sistema debe permitir al vendedor ver los pedidos asociados a
                  su tienda, con detalle de productos y datos del comprador. \\
  \rowcolor{rowalt}
  Prioridad     & Media \\
  Rol afectado  & Vendedor \\
  \bottomrule
\end{tabularx}

% ── RF-10 ────────────────────────────────────────────────────────────────────
\vspace{6pt}
\begin{tabularx}{\linewidth}{>{\bfseries\color{primary}}p{3.5cm}L}
  \toprule
  \rowcolor{headerbg}
  \multicolumn{2}{l}{\textcolor{white}{\textbf{RF-10 — Panel de administración}}} \\
  \midrule
  Descripción   & El sistema debe proporcionar al administrador herramientas para
                  gestionar usuarios, aprobar o suspender tiendas y moderar productos. \\
  \rowcolor{rowalt}
  Prioridad     & Media \\
  Rol afectado  & Administrador \\
  \bottomrule
\end{tabularx}

% ── RF-11 ────────────────────────────────────────────────────────────────────
\vspace{6pt}
\begin{tabularx}{\linewidth}{>{\bfseries\color{primary}}p{3.5cm}L}
  \toprule
  \rowcolor{headerbg}
  \multicolumn{2}{l}{\textcolor{white}{\textbf{RF-11 — Gestión de perfil}}} \\
  \midrule
  Descripción   & El sistema debe permitir a cualquier usuario autenticado
                  visualizar y editar su información personal. \\
  \rowcolor{rowalt}
  Prioridad     & Baja \\
  Rol afectado  & Cliente, Vendedor \\
  \bottomrule
\end{tabularx}

% ── RF-12 ────────────────────────────────────────────────────────────────────
\vspace{6pt}
\begin{tabularx}{\linewidth}{>{\bfseries\color{primary}}p{3.5cm}L}
  \toprule
  \rowcolor{headerbg}
  \multicolumn{2}{l}{\textcolor{white}{\textbf{RF-12 — Cerrar sesión}}} \\
  \midrule
  Descripción   & El sistema debe permitir al usuario cerrar su sesión de forma
                  segura, invalidando la sesión activa. \\
  \rowcolor{rowalt}
  Prioridad     & Alta \\
  Rol afectado  & Cliente, Vendedor, Administrador \\
  \bottomrule
\end{tabularx}

% ════════════════════════════════════════════════════════════════════════════
\section{Requerimientos No Funcionales}
% ════════════════════════════════════════════════════════════════════════════

\begin{tabularx}{\linewidth}{>{\bfseries\color{primary}}p{1.5cm}>{\bfseries}p{3cm}X}
  \toprule
  \rowcolor{headerbg}
  \textcolor{white}{\textbf{ID}} &
  \textcolor{white}{\textbf{Categoría}} &
  \textcolor{white}{\textbf{Descripción}} \\
  \midrule
  RNF-01 & Usabilidad
    & El sistema debe ser intuitivo y fácil de usar para personas con conocimientos
      básicos de navegación web. \\
  \rowcolor{rowalt}
  RNF-02 & Seguridad
    & El sistema debe requerir autenticación para acceder a cualquier funcionalidad
      protegida. Las contraseñas deben almacenarse con \emph{hashing}. \\
  RNF-03 & Rendimiento
    & El sistema debe responder a las acciones del usuario en menos de 2 segundos
      bajo condiciones normales de uso. \\
  \rowcolor{rowalt}
  RNF-04 & Compatibilidad
    & El sistema debe ser accesible desde navegadores web modernos (Chrome, Firefox,
      Edge, Safari) sin requerir plugins adicionales. \\
  RNF-05 & Disponibilidad
    & El sistema debe estar disponible al menos el 95\% del tiempo durante el
      periodo de evaluación. \\
  \rowcolor{rowalt}
  RNF-06 & Mantenibilidad
    & El código debe seguir las convenciones de Django y estar organizado en
      aplicaciones independientes para facilitar su extensión. \\
  \bottomrule
\end{tabularx}

\newpage
% ════════════════════════════════════════════════════════════════════════════
\section{Casos de Uso}
% ════════════════════════════════════════════════════════════════════════════

% ─────────────────────────────────────────────────────────────────────────────
\subsection{CU-01: Registrarse en la plataforma}

\begin{tabularx}{\linewidth}{>{\bfseries\color{primary}}p{3.5cm}X}
  \toprule
  Actor principal & Usuario no registrado \\
  \rowcolor{rowalt}
  Precondición    & El usuario no tiene cuenta en la plataforma \\
  \midrule
  \rowcolor{primarylight}
  \multicolumn{2}{l}{\textbf{Flujo principal}} \\
  \midrule
  Paso 1 & El usuario accede a la página de registro \\
  \rowcolor{rowalt}
  Paso 2 & Ingresa nombre, correo electrónico y contraseña \\
  Paso 3 & El sistema valida que el correo no esté en uso y que la contraseña cumpla los requisitos \\
  \rowcolor{rowalt}
  Paso 4 & El sistema crea la cuenta y redirige al usuario al inicio \\
  \midrule
  \rowcolor{primarylight}
  \multicolumn{2}{l}{\textbf{Flujos alternativos}} \\
  \midrule
  A1 & Correo ya registrado → el sistema muestra un mensaje de error \\
  \rowcolor{rowalt}
  A2 & Contraseña inválida → el sistema indica los requisitos no cumplidos \\
  \bottomrule
\end{tabularx}

% ─────────────────────────────────────────────────────────────────────────────
\subsection{CU-02: Iniciar sesión}

\begin{tabularx}{\linewidth}{>{\bfseries\color{primary}}p{3.5cm}X}
  \toprule
  Actor principal & Usuario registrado \\
  \rowcolor{rowalt}
  Precondición    & El usuario tiene una cuenta activa \\
  \midrule
  \rowcolor{primarylight}
  \multicolumn{2}{l}{\textbf{Flujo principal}} \\
  \midrule
  Paso 1 & El usuario accede a la página de inicio de sesión \\
  \rowcolor{rowalt}
  Paso 2 & Ingresa su correo electrónico y contraseña \\
  Paso 3 & El sistema valida las credenciales \\
  \rowcolor{rowalt}
  Paso 4 & El sistema inicia la sesión y redirige según el rol del usuario \\
  \midrule
  \rowcolor{primarylight}
  \multicolumn{2}{l}{\textbf{Flujos alternativos}} \\
  \midrule
  A1 & Credenciales incorrectas → el sistema muestra un mensaje de error sin revelar cuál campo es incorrecto \\
  \bottomrule
\end{tabularx}

% ─────────────────────────────────────────────────────────────────────────────
\subsection{CU-03: Explorar y buscar productos}

\begin{tabularx}{\linewidth}{>{\bfseries\color{primary}}p{3.5cm}X}
  \toprule
  Actor principal & Cliente \\
  \midrule
  \rowcolor{primarylight}
  \multicolumn{2}{l}{\textbf{Flujo principal}} \\
  \midrule
  Paso 1 & El cliente navega al catálogo de productos \\
  \rowcolor{rowalt}
  Paso 2 & El sistema muestra todos los productos disponibles \\
  Paso 3 & El cliente aplica un filtro por nombre o categoría \\
  \rowcolor{rowalt}
  Paso 4 & El sistema muestra los resultados filtrados \\
  \midrule
  \rowcolor{primarylight}
  \multicolumn{2}{l}{\textbf{Flujos alternativos}} \\
  \midrule
  A1 & Sin resultados → el sistema muestra un mensaje indicando que no se encontraron productos \\
  \bottomrule
\end{tabularx}

% ─────────────────────────────────────────────────────────────────────────────
\subsection{CU-04: Agregar producto al carrito}

\begin{tabularx}{\linewidth}{>{\bfseries\color{primary}}p{3.5cm}X}
  \toprule
  Actor principal & Cliente \\
  \rowcolor{rowalt}
  Precondición    & El cliente está autenticado y el producto tiene stock disponible \\
  \midrule
  \rowcolor{primarylight}
  \multicolumn{2}{l}{\textbf{Flujo principal}} \\
  \midrule
  Paso 1 & El cliente selecciona un producto del catálogo \\
  \rowcolor{rowalt}
  Paso 2 & Visualiza el detalle del producto \\
  Paso 3 & Selecciona la cantidad y hace clic en «Añadir al carrito» \\
  \rowcolor{rowalt}
  Paso 4 & El sistema agrega el producto al carrito y actualiza el contador \\
  \midrule
  \rowcolor{primarylight}
  \multicolumn{2}{l}{\textbf{Flujos alternativos}} \\
  \midrule
  A1 & Producto sin stock → el sistema deshabilita el botón e informa la situación \\
  \bottomrule
\end{tabularx}

% ─────────────────────────────────────────────────────────────────────────────
\subsection{CU-05: Realizar un pedido}

\begin{tabularx}{\linewidth}{>{\bfseries\color{primary}}p{3.5cm}X}
  \toprule
  Actor principal & Cliente \\
  \rowcolor{rowalt}
  Precondición    & El carrito tiene al menos un producto \\
  \midrule
  \rowcolor{primarylight}
  \multicolumn{2}{l}{\textbf{Flujo principal}} \\
  \midrule
  Paso 1 & El cliente accede a su carrito \\
  \rowcolor{rowalt}
  Paso 2 & Revisa los productos y cantidades \\
  Paso 3 & Confirma el pedido \\
  \rowcolor{rowalt}
  Paso 4 & El sistema registra el pedido, actualiza el inventario y muestra confirmación \\
  \midrule
  \rowcolor{primarylight}
  \multicolumn{2}{l}{\textbf{Flujos alternativos}} \\
  \midrule
  A1 & Algún producto ya no tiene stock suficiente → el sistema informa al cliente antes de confirmar \\
  \bottomrule
\end{tabularx}

% ─────────────────────────────────────────────────────────────────────────────
\subsection{CU-06: Crear tienda}

\begin{tabularx}{\linewidth}{>{\bfseries\color{primary}}p{3.5cm}X}
  \toprule
  Actor principal & Vendedor \\
  \rowcolor{rowalt}
  Precondición    & El usuario está autenticado y no tiene una tienda creada \\
  \midrule
  \rowcolor{primarylight}
  \multicolumn{2}{l}{\textbf{Flujo principal}} \\
  \midrule
  Paso 1 & El vendedor accede a la opción «Crear tienda» \\
  \rowcolor{rowalt}
  Paso 2 & Ingresa nombre, descripción y categoría principal \\
  Paso 3 & El sistema registra la tienda con estado \textit{pendiente} \\
  \rowcolor{rowalt}
  Paso 4 & El administrador aprueba la tienda \\
  Paso 5 & La tienda queda activa y visible en la plataforma \\
  \midrule
  \rowcolor{primarylight}
  \multicolumn{2}{l}{\textbf{Flujos alternativos}} \\
  \midrule
  A1 & Nombre de tienda ya en uso → el sistema solicita un nombre diferente \\
  \rowcolor{rowalt}
  A2 & Administrador rechaza la tienda → el vendedor recibe una notificación \\
  \bottomrule
\end{tabularx}

% ─────────────────────────────────────────────────────────────────────────────
\subsection{CU-07: Publicar producto}

\begin{tabularx}{\linewidth}{>{\bfseries\color{primary}}p{3.5cm}X}
  \toprule
  Actor principal & Vendedor \\
  \rowcolor{rowalt}
  Precondición    & El vendedor tiene una tienda activa \\
  \midrule
  \rowcolor{primarylight}
  \multicolumn{2}{l}{\textbf{Flujo principal}} \\
  \midrule
  Paso 1 & El vendedor accede al panel de su tienda \\
  \rowcolor{rowalt}
  Paso 2 & Selecciona «Agregar producto» \\
  Paso 3 & Ingresa nombre, descripción, precio, cantidad e imagen \\
  \rowcolor{rowalt}
  Paso 4 & El sistema registra el producto y lo hace visible en el catálogo \\
  \midrule
  \rowcolor{primarylight}
  \multicolumn{2}{l}{\textbf{Flujos alternativos}} \\
  \midrule
  A1 & Datos incompletos → el sistema indica los campos faltantes \\
  \bottomrule
\end{tabularx}

% ─────────────────────────────────────────────────────────────────────────────
\subsection{CU-08: Administrar la plataforma}

\begin{tabularx}{\linewidth}{>{\bfseries\color{primary}}p{3.5cm}X}
  \toprule
  Actor principal & Administrador \\
  \rowcolor{rowalt}
  Precondición    & El administrador está autenticado con privilegios de superusuario \\
  \midrule
  \rowcolor{primarylight}
  \multicolumn{2}{l}{\textbf{Flujo principal}} \\
  \midrule
  Paso 1 & El administrador accede al panel de administración \\
  \rowcolor{rowalt}
  Paso 2 & Selecciona la entidad a gestionar (usuarios, tiendas o productos) \\
  Paso 3 & Realiza la acción: aprobar, suspender o eliminar \\
  \rowcolor{rowalt}
  Paso 4 & El sistema aplica el cambio y actualiza la vista \\
  \bottomrule
\end{tabularx}

\newpage
% ════════════════════════════════════════════════════════════════════════════
\section{Restricciones}
% ════════════════════════════════════════════════════════════════════════════

\begin{tabularx}{\linewidth}{>{\bfseries\color{primary}}p{3.5cm}X}
  \toprule
  \rowcolor{primarylight}
  Restricción & Descripción \\
  \midrule
  Tecnología backend
    & El sistema debe desarrollarse obligatoriamente con el framework Django
      (Python). \\
  Frontend
    & La interfaz se construirá con HTML, CSS y JavaScript estándar, sin frameworks
      de frontend de gran escala. \\
  Pagos
    & No se integrará una pasarela de pago real en esta versión del proyecto. \\
  Tiempo
    & El desarrollo se realizará en el periodo estipulado por el curso; los
      requerimientos deben ser estables para respetar el cronograma. \\
  \bottomrule
\end{tabularx}

% ════════════════════════════════════════════════════════════════════════════
\section{Mapa de Navegación}
% ════════════════════════════════════════════════════════════════════════════
\begin{center}
    \includegraphics[width=1\textwidth]{Mapa de Navegación.png}
\end{center}

\newpage
% ════════════════════════════════════════════════════════════════════════════
\section{Caso de Uso}
% ════════════════════════════════════════════════════════════════════════════

\begin{center}
    \includegraphics[width=0.60\textwidth]{Diagramas de actividad/caso de uso.png}
\end{center}
\newpage
% ════════════════════════════════════════════════════════════════════════════
\section{Diagramas de Actividad}
% ════════════════════════════════════════════════════════════════════════════

\begin{figure}[H]
    \centering
    \begin{subfigure}{0.40\textwidth}
        \includegraphics[width=\textwidth]{Diagramas de actividad/Registro de cuentas bueno.png}
        \caption{Registro de cuenta}
    \end{subfigure}
    \hfill
    \begin{subfigure}{0.40\textwidth}
        \includegraphics[width=\textwidth]{Diagramas de actividad/Inicio de sesión.png}
        \caption{Inicio de sesión}
    \end{subfigure}

    \begin{subfigure}{0.40\textwidth}
        \includegraphics[width=\textwidth]{Diagramas de actividad/Buscar productos.png}
        \caption{Buscar productos}
    \end{subfigure}
    \hfill
    \begin{subfigure}{0.40\textwidth}
        \includegraphics[width=\textwidth]{Diagramas de actividad/Agregar al carrito.png}
        \caption{Agregar al carrito}
    \end{subfigure}
\end{figure}

\begin{figure}[H]
    \centering
    \begin{subfigure}{0.40\textwidth}
        \includegraphics[width=\textwidth]{Diagramas de actividad/Realizar pedido.png}
        \caption{Realizar pedido}
    \end{subfigure}
    \hfill
    \begin{subfigure}{0.40\textwidth}
        \includegraphics[width=\textwidth]{Diagramas de actividad/Crear Tienda.png}
        \caption{Crear tienda}
    \end{subfigure}


    \begin{subfigure}{0.40\textwidth}
        \includegraphics[width=\textwidth]{Diagramas de actividad/Publicar Producto.png}
        \caption{Publicar producto}
    \end{subfigure}
    \hfill
    \begin{subfigure}{0.40\textwidth}
        \includegraphics[width=\textwidth]{Diagramas de actividad/Administrar plataforma.png}
        \caption{Administrar plataforma}
    \end{subfigure}
\end{figure}
\clearpage
\end{document}
